\documentclass[11pt,a4paper]{article}

\usepackage[utf8]{inputenc}
\usepackage[T1]{fontenc}
\usepackage[french]{babel}
\usepackage[a4paper,margin=1.7cm]{geometry}
\usepackage{amsmath,amssymb,amsthm,mathtools}
\usepackage{lmodern}
\usepackage{xcolor}
\usepackage{array,tabularx,booktabs,longtable,multirow}
\usepackage{enumitem}
\usepackage{fancyhdr}
\usepackage{lastpage}
\usepackage{hyperref}
\usepackage{graphicx}
\usepackage{tikz}
\usepackage{pgfplots}
\usepackage[most]{tcolorbox}
\usepackage{multicol}
\usepackage{siunitx}
\usepackage{stmaryrd}
\usetikzlibrary{arrows.meta,calc,positioning,patterns,decorations.pathmorphing}
\pgfplotsset{compat=1.18}
\hypersetup{colorlinks=true,linkcolor=blue,urlcolor=blue}
\sisetup{locale=FR}

\setlength{\headheight}{14pt}
\pagestyle{fancy}
\fancyhf{}
\fancyhead[L]{Première générale -- Spécialité Physique-Chimie}
\fancyhead[C]{Fiche d'exercices -- Ondes et signaux}
\fancyhead[R]{Page \thepage/\pageref{LastPage}}
\fancyfoot[C]{Pierre Garcia-Buscaylet -- fiche progressive et approfondie}

\definecolor{bleu1}{RGB}{30,90,160}
\definecolor{bleu2}{RGB}{232,240,255}
\definecolor{vert1}{RGB}{30,120,80}
\definecolor{vert2}{RGB}{232,247,239}
\definecolor{rouge1}{RGB}{165,55,55}
\definecolor{rouge2}{RGB}{252,238,238}
\definecolor{jaune2}{RGB}{255,248,220}
\definecolor{violet2}{RGB}{244,238,255}
\definecolor{grisclair}{RGB}{246,247,249}

\newtcolorbox{objectifbox}{colback=bleu2,colframe=bleu1,title=Objectif,fonttitle=\bfseries}
\newtcolorbox{rappelbox}{colback=vert2,colframe=vert1,title=Rappel utile,fonttitle=\bfseries}
\newtcolorbox{methodebox}{colback=jaune2,colframe=orange!80!black,title=Méthode générale,fonttitle=\bfseries}
\newtcolorbox{piegebox}{colback=rouge2,colframe=rouge1,title=Piège classique,fonttitle=\bfseries}
\newtcolorbox{ouverturebox}{colback=violet2,colframe=purple!70!black,title=Ouverture Terminale,fonttitle=\bfseries}
\newtcolorbox{baremebox}{colback=grisclair,colframe=black!70,title=Repères de difficulté,fonttitle=\bfseries}
\newtcolorbox{corrbox}{colback=white,colframe=black!75,title=Éléments de correction,fonttitle=\bfseries}

\newcommand{\etoiles}[1]{%
\ifcase#1\or $\star$\or $\star\star$\or $\star\star\star$\or $\star\star\star\star$\or $\star\star\star\star\star$\fi}
\newcommand{\exostar}[2]{\subsection*{Exercice #1 \hfill \etoiles{#2}}}
\newcommand{\vrai}{\textbf{Vrai}}
\newcommand{\faux}{\textbf{Faux}}

\begin{document}

\begin{center}
    {\LARGE\bfseries Grosse fiche d'exercices -- Ondes et signaux}\\[0.3em]
    {\large Première spécialité Physique-Chimie}\\[0.4em]
    \rule{0.92\linewidth}{0.6pt}
\end{center}

\begin{objectifbox}
Cette fiche est conçue pour \textbf{faire comprendre} et pas seulement pour appliquer des formules. Elle mélange :
\begin{itemize}[leftmargin=1.2em]
    \item des \textbf{QCM} pour tester les réflexes ;
    \item des \textbf{vrai/faux argumentés} pour traquer les confusions ;
    \item des \textbf{applications directes} pour consolider les méthodes ;
    \item des \textbf{problèmes longs et progressifs} ;
    \item des \textbf{problèmes plus ouverts et plus exigeants} pour sortir l'élève de sa zone de confort.
\end{itemize}
Les exercices sont classés avec un \textbf{gradient de difficulté} de $\star$ à $\star\star\star\star\star$.
\end{objectifbox}

\begin{baremebox}
\begin{itemize}[leftmargin=1.2em]
    \item $\star$ : application immédiate du cours ;
    \item $\star\star$ : exercice classique avec un ou deux pièges ;
    \item $\star\star\star$ : raisonnement structuré ;
    \item $\star\star\star\star$ : problème long ou moins guidé ;
    \item $\star\star\star\star\star$ : problème difficile, exigeant sur le sens physique et l'autonomie.
\end{itemize}
\end{baremebox}

\tableofcontents
\bigskip

\section{QCM de démarrage}

\exostar{1}{1}
Pour chaque proposition, choisir la bonne réponse. Une seule réponse est attendue sauf mention contraire.

\begin{enumerate}[label=\textbf{Q\arabic*.},leftmargin=1.1cm]
\item Une onde mécanique progressive est :
\begin{enumerate}[label=\alph*)]
    \item un transport de matière sans transport d'énergie ;
    \item un transport d'énergie sans transport global de matière ;
    \item une vibration qui ne se propage pas ;
    \item un rayonnement électromagnétique.
\end{enumerate}

\item Le son ne se propage pas dans le vide car :
\begin{enumerate}[label=\alph*)]
    \item sa fréquence y devient nulle ;
    \item sa longueur d'onde y devient infinie ;
    \item il a besoin d'un milieu matériel ;
    \item l'air absorbe toujours le son.
\end{enumerate}

\item Une onde se propage à la célérité \SI{4,0}{m.s^{-1}}. En \SI{0,50}{s}, elle parcourt :
\begin{enumerate}[label=\alph*)]
    \item \SI{0,125}{m}
    \item \SI{2,0}{m}
    \item \SI{4,5}{m}
    \item \SI{8,0}{m}
\end{enumerate}

\item Pour une onde périodique, la relation correcte est :
\begin{enumerate}[label=\alph*)]
    \item $\lambda = \dfrac{T}{c}$
    \item $c = \lambda T$
    \item $\lambda = cT$
    \item $T = c f$
\end{enumerate}

\item Deux points d'un milieu séparés d'une distance égale à une longueur d'onde vibrent :
\begin{enumerate}[label=\alph*)]
    \item en opposition de phase ;
    \item en phase ;
    \item avec des amplitudes différentes ;
    \item avec des fréquences différentes.
\end{enumerate}

\item La couleur perçue d'un objet dépend :
\begin{enumerate}[label=\alph*)]
    \item seulement de la couleur propre de l'objet ;
    \item seulement de la lumière incidente ;
    \item de la lumière incidente et des phénomènes d'absorption, transmission et diffusion ;
    \item uniquement de la présence d'un filtre.
\end{enumerate}

\item Une lentille mince convergente peut former d'un objet réel :
\begin{enumerate}[label=\alph*)]
    \item seulement des images réelles ;
    \item seulement des images virtuelles ;
    \item des images réelles ou virtuelles selon la position de l'objet ;
    \item jamais d'image droite.
\end{enumerate}

\item L'énergie d'un photon est donnée par :
\begin{enumerate}[label=\alph*)]
    \item $E=h/\nu$
    \item $E=h\nu$
    \item $E=c\lambda$
    \item $E=hc\lambda$
\end{enumerate}

\item Quand la fréquence de la lumière augmente :
\begin{enumerate}[label=\alph*)]
    \item l'énergie du photon diminue ;
    \item la longueur d'onde augmente forcément ;
    \item l'énergie du photon augmente ;
    \item la célérité dans le vide change.
\end{enumerate}

\item Dans le vide, on peut écrire :
\begin{enumerate}[label=\alph*)]
    \item $c=\lambda\nu$
    \item $c=\lambda/\nu$
    \item $\nu = c\lambda$
    \item $E=mc^2$ pour décrire un photon au programme de première.
\end{enumerate}
\end{enumerate}

\section{Vrai/faux pour comprendre finement}

\exostar{2}{2}
Indiquer si chaque affirmation est vraie ou fausse, puis \textbf{justifier en une ou deux phrases précises}. Une réponse non justifiée ne compte pas.

\begin{enumerate}[label=\textbf{Affirmation \arabic*.},leftmargin=1.8cm]
    \item Une onde transporte toujours de la matière du point source vers le récepteur.
    \item Une onde mécanique peut se propager dans le vide si elle a une grande amplitude.
    \item La célérité d'une onde dépend du milieu de propagation.
    \item La fréquence d'une onde est imposée par la source.
    \item Si la célérité change lorsqu'une onde passe dans un autre milieu, sa fréquence change nécessairement aussi.
    \item La longueur d'onde dépend à la fois de la source et du milieu.
    \item Deux points séparés de $\lambda/2$ vibrent forcément en phase.
    \item Une image réelle formée par une lentille convergente est nécessairement renversée.
    \item Un objet rouge éclairé en lumière bleue apparaît rouge.
    \item Plus la longueur d'onde d'un photon est petite, plus son énergie est grande.
    \item Un spectre de raies d'émission est lié à la quantification des niveaux d'énergie atomiques.
    \item Le modèle ondulatoire et le modèle particulaire de la lumière s'excluent absolument.
\end{enumerate}

\section{Applications directes et méthodes fondamentales}

\begin{rappelbox}
Relations à connaître sans hésitation :
\[
 c = \frac{d}{\Delta t},\qquad \lambda = cT,\qquad f=\frac{1}{T},\qquad c=\lambda f,\qquad E=h\nu = \frac{hc}{\lambda}.
\]
En optique géométrique, avec les conventions algébriques usuelles :
\[
\frac{1}{\overline{OA'}}-\frac{1}{\overline{OA}}=\frac{1}{f'},
\qquad \gamma = \frac{\overline{A'B'}}{\overline{AB}} = \frac{\overline{OA'}}{\overline{OA}}.
\]
\end{rappelbox}

\exostar{3}{1}
Sur une corde, une impulsion se déplace à la célérité \SI{5,0}{m.s^{-1}}. Deux points $A$ et $B$ sont séparés de \SI{1,5}{m}.
\begin{enumerate}[label=\alph*)]
    \item Calculer le retard de l'onde entre $A$ et $B$.
    \item Interpréter physiquement ce retard.
\end{enumerate}

\exostar{4}{1}
Une onde sonore se propage dans l'air à \SI{340}{m.s^{-1}}. Un écho est entendu \SI{0,80}{s} après l'émission d'un clap.
\begin{enumerate}[label=\alph*)]
    \item Quelle distance le son a-t-il parcourue au total ?
    \item À quelle distance se trouve l'obstacle ?
    \item Pourquoi faut-il diviser par 2 à la fin ?
\end{enumerate}

\exostar{5}{2}
Une onde périodique de fréquence \SI{25}{Hz} se propage à la célérité \SI{3,0}{m.s^{-1}} le long d'un ressort.
\begin{enumerate}[label=\alph*)]
    \item Déterminer la période.
    \item Déterminer la longueur d'onde.
    \item Donner la distance entre deux points vibrant en phase.
    \item Donner la distance entre deux points en opposition de phase.
\end{enumerate}

\exostar{6}{2}
Une représentation temporelle d'un signal montre 4 périodes sur une durée de \SI{20}{ms}.
\begin{enumerate}[label=\alph*)]
    \item Déterminer la période $T$.
    \item Déterminer la fréquence $f$.
    \item Le signal serait-il plus aigu ou plus grave si sa période était doublée ? Expliquer dans le cas d'un son.
\end{enumerate}

\exostar{7}{2}
Sur une photographie stroboscopique d'une onde à la surface de l'eau, la distance entre deux crêtes successives vaut \SI{12}{cm}. La période de la source est \SI{0,40}{s}.
\begin{enumerate}[label=\alph*)]
    \item Identifier la grandeur \SI{12}{cm}.
    \item Calculer la célérité de l'onde.
    \item Que devient la longueur d'onde si on double la fréquence de la source sans changer le milieu ?
\end{enumerate}

\section{Optique : images, lentilles et couleurs}

\exostar{8}{2}
Un objet de hauteur \SI{2,0}{cm} est placé à \SI{30}{cm} devant une lentille convergente de distance focale \SI{10}{cm}.
\begin{enumerate}[label=\alph*)]
    \item Déterminer la position de l'image à l'aide de la relation de conjugaison.
    \item Calculer le grandissement.
    \item En déduire la taille de l'image.
    \item Préciser si l'image est réelle ou virtuelle, droite ou renversée.
\end{enumerate}

\exostar{9}{3}
On veut projeter nettement sur un écran l'image d'un objet lumineux à l'aide d'une lentille convergente.
\begin{enumerate}[label=\alph*)]
    \item Pourquoi l'image obtenue sur l'écran est-elle nécessairement réelle ?
    \item L'objet est placé entre le foyer et la lentille. Peut-on encore obtenir une image nette sur l'écran ? Justifier.
    \item Citer deux moyens physiques de réaliser la mise au point dans un système optique réel.
    \item Expliquer qualitativement pourquoi les appareils photo et l'œil doivent faire la mise au point.
\end{enumerate}

\exostar{10}{2}
On éclaire successivement un carton cyan, un carton magenta et un carton jaune par différentes lumières.
\begin{enumerate}[label=\alph*)]
    \item Quelles composantes un objet cyan diffuse-t-il principalement ?
    \item Quelle couleur perçoit-on en éclairant un objet jaune avec une lumière rouge ?
    \item Quelle couleur perçoit-on en éclairant le même objet jaune avec une lumière bleue ?
    \item Expliquer la différence entre synthèse additive et synthèse soustractive.
\end{enumerate}

\section{Photons, spectres et niveaux d'énergie}

\exostar{11}{2}
Calculer l'énergie d'un photon de longueur d'onde \SI{650}{nm}. On prendra
\[
 h = 6{,}63\times 10^{-34}\ \si{J.s},\qquad c=3{,}00\times10^8\ \si{m.s^{-1}}.
\]
Donner le résultat en joules puis commenter qualitativement s'il s'agit d'une lumière plutôt énergétique ou non au sein du visible.

\exostar{12}{3}
On considère un atome possédant deux niveaux d'énergie $E_1$ et $E_2$ tels que
\[
E_2-E_1 = 3{,}0\times 10^{-19}\ \si{J}.
\]
\begin{enumerate}[label=\alph*)]
    \item Calculer la fréquence du photon émis lors de la transition de $E_2$ vers $E_1$.
    \item En déduire sa longueur d'onde dans le vide.
    \item Situer qualitativement cette radiation : ultraviolet, visible ou infrarouge si l'on rappelle que le visible va approximativement de \SI{400}{nm} à \SI{800}{nm}.
    \item Expliquer pourquoi on obtient des raies et non un spectre continu pour une source atomique diluée.
\end{enumerate}

\section{Problèmes progressifs}

\exostar{13}{3}
\textbf{Repérage d'un orage.} Un observateur voit un éclair puis entend le tonnerre \SI{5,4}{s} plus tard.
\begin{enumerate}[label=\alph*)]
    \item Pourquoi peut-on considérer que l'éclair est perçu quasiment instantanément à l'échelle de l'expérience ?
    \item En prenant la célérité du son dans l'air égale à \SI{340}{m.s^{-1}}, estimer la distance de l'orage.
    \item Un second observateur situé \SI{680}{m} plus loin dans la direction de l'orage entend le tonnerre avec quel décalage par rapport au premier ?
    \item Expliquer ce que cette situation révèle sur la différence de célérité entre onde lumineuse et onde sonore.
\end{enumerate}

\exostar{14}{3}
\textbf{Onde le long d'un ressort.} On filme une onde périodique se propageant le long d'un ressort horizontal. Sur les images, on mesure une distance de \SI{45}{cm} entre trois crêtes consécutives. La source effectue 10 allers-retours complets en \SI{4,0}{s}.
\begin{enumerate}[label=\alph*)]
    \item Déterminer la période de la source.
    \item Déterminer la fréquence.
    \item Déterminer la longueur d'onde.
    \item En déduire la célérité.
    \item Un point situé à \SI{90}{cm} de la source commence-t-il à osciller avant ou après \SI{0,50}{s} ? Justifier.
\end{enumerate}

\exostar{15}{3}
\textbf{Mise au point et taille d'image.} On photographie une petite figurine de hauteur \SI{12}{cm}. L'objectif est modélisé par une lentille convergente de focale \SI{50}{mm}. Le capteur se trouve à \SI{60}{mm} derrière la lentille quand l'image est nette.
\begin{enumerate}[label=\alph*)]
    \item Déterminer la distance objet-lentille.
    \item Calculer le grandissement.
    \item En déduire la taille de l'image sur le capteur.
    \item L'image est-elle droite ou renversée ?
    \item Pourquoi une augmentation de la distance objet-lentille oblige-t-elle à modifier la position du capteur ou la focale pour conserver une image nette ?
\end{enumerate}

\exostar{16}{4}
\textbf{Couleur d'un objet selon l'éclairage.} On considère trois objets idéalisés : un objet rouge, un objet cyan et un objet blanc. On les éclaire successivement en lumière blanche, rouge, verte puis bleue.
\begin{enumerate}[label=\alph*)]
    \item Compléter un tableau donnant la couleur apparente de chaque objet pour chaque éclairage.
    \item Expliquer en détail ce que diffuse ou absorbe chaque objet.
    \item Pourquoi un objet blanc apparaît-il de la couleur de la lumière qui l'éclaire ?
    \item Un objet cyan peut-il apparaître noir ? Donner une situation précise.
    \item En quoi cet exercice aide-t-il à comprendre les erreurs fréquentes sur la ``couleur propre'' d'un objet ?
\end{enumerate}

\exostar{17}{4}
\textbf{Spectre et transitions atomiques.} Un gaz excité émet trois raies de longueurs d'onde \SI{410}{nm}, \SI{434}{nm} et \SI{486}{nm}.
\begin{enumerate}[label=\alph*)]
    \item Classer ces radiations de la plus énergétique à la moins énergétique.
    \item Calculer l'énergie des photons associés à chacune des trois raies.
    \item Expliquer pourquoi l'observation de quelques raies bien définies plaide en faveur d'une quantification des niveaux d'énergie.
    \item Une radiation de \SI{820}{nm} appartiendrait-elle encore au visible ?
    \item Quel changement de niveau d'énergie est le plus grand parmi les trois transitions considérées ?
\end{enumerate}

\section{Problèmes longs pour sortir de la zone de confort}

\exostar{18}{4}
\textbf{Localisation d'une source sonore par deux microphones.}
Deux microphones $M_1$ et $M_2$ sont disposés sur une même ligne, séparés de \SI{3,40}{m}. Une source sonore ponctuelle se trouve quelque part sur cette ligne. On observe que le signal arrive au microphone $M_1$ avec une avance de \SI{4,0}{ms} sur le signal reçu par $M_2$. On prendra la célérité du son égale à \SI{340}{m.s^{-1}}.
\begin{enumerate}[label=\alph*)]
    \item Convertir l'avance en seconde.
    \item À quelle différence de distance source-micro correspond ce retard ?
    \item En déduire où peut se trouver la source sur la ligne reliant les deux microphones.
    \item Déterminer sa position précise si elle est située entre les deux microphones.
    \item Déterminer sa position précise si elle est située à l'extérieur du segment $[M_1M_2]$ du côté de $M_1$.
    \item Discuter : la seule mesure de retard suffit-elle toujours à localiser une source en 2D ou en 3D ?
\end{enumerate}

\exostar{19}{4}
\textbf{Lentille et changement de focale.}
On modélise un appareil photo par une lentille mince convergente. On photographie un objet placé à \SI{2,0}{m}. Le capteur est fixe à \SI{52}{mm} derrière la lentille.
\begin{enumerate}[label=\alph*)]
    \item Quelle focale faut-il choisir pour obtenir une image nette sur le capteur ?
    \item On remplace l'objectif par un autre de focale plus grande, sans déplacer le capteur. L'image devient-elle nette ?
    \item Le grandissement augmente-t-il ou diminue-t-il quand la focale augmente dans cette situation ?
    \item Expliquer le lien qualitatif avec l'idée de zoom.
    \item Pourquoi un zoom optique n'est-il pas seulement un ``agrandissement informatique'' ?
\end{enumerate}

\exostar{20}{5}
\textbf{Problème de synthèse : un capteur scientifique observe une raie lumineuse.}
Un capteur détecte une raie de longueur d'onde \SI{486}{nm} émise par un gaz excité. L'instrument d'observation comporte une lentille convergente de distance focale \SI{100}{mm}. L'objet lumineux observé est à \SI{1,20}{m} de la lentille.
\begin{enumerate}[label=\alph*)]
    \item Déterminer la position de l'image formée par la lentille.
    \item Déterminer le grandissement si la taille de l'objet est \SI{3,0}{mm}.
    \item Déterminer la taille de l'image.
    \item Calculer l'énergie du photon associé à la raie de \SI{486}{nm}.
    \item Expliquer pourquoi une raie spectrale permet d'identifier une espèce chimique.
    \item Expliquer en quoi ce problème mélange les deux modèles de la lumière : géométrique/ondulatoire d'un côté, particulaire de l'autre.
    \item Proposer deux erreurs fréquentes qu'un élève pourrait commettre dans ce problème et expliquer comment les éviter.
\end{enumerate}

\exostar{21}{5}
\textbf{Problème ouvert : peut-on tout déduire d'une seule mesure ?}
Une source périodique crée une onde mécanique sinusoïdale à la surface de l'eau. Un élève mesure uniquement la période temporelle en un point et trouve $T=\SI{0,20}{s}$.
\begin{enumerate}[label=\alph*)]
    \item Peut-il en déduire la longueur d'onde ?
    \item Peut-il en déduire la célérité ?
    \item Quelles informations supplémentaires minimales faut-il obtenir expérimentalement pour calculer toutes les autres grandeurs ?
    \item Proposer deux protocoles expérimentaux différents permettant de déterminer la célérité.
    \item Expliquer pourquoi cet exercice, apparemment simple, est en réalité un test très profond de compréhension des grandeurs d'une onde.
\end{enumerate}

\section{QCM d'approfondissement}

\exostar{22}{3}
Choisir la ou les bonnes réponses.
\begin{enumerate}[label=\textbf{QCM \arabic*.},leftmargin=1.5cm]
    \item Pour une onde périodique se propageant dans un milieu donné :
    \begin{enumerate}[label=\alph*)]
        \item si la fréquence double, la longueur d'onde double ;
        \item si la fréquence double, la longueur d'onde est divisée par deux ;
        \item si l'amplitude double, la célérité double ;
        \item la célérité dépend uniquement de la fréquence.
    \end{enumerate}
    \item Un objet réel placé entre une lentille convergente et son foyer donne :
    \begin{enumerate}[label=\alph*)]
        \item une image réelle sur un écran ;
        \item une image virtuelle, droite, agrandie ;
        \item une image virtuelle, renversée ;
        \item aucune image du tout.
    \end{enumerate}
    \item Un photon ultraviolet, par rapport à un photon rouge visible :
    \begin{enumerate}[label=\alph*)]
        \item a une plus grande fréquence ;
        \item a une plus petite énergie ;
        \item a une plus petite longueur d'onde ;
        \item se propage plus vite dans le vide.
    \end{enumerate}
    \item À propos d'un spectre de raies :
    \begin{enumerate}[label=\alph*)]
        \item il correspond à toutes les longueurs d'onde possibles ;
        \item il traduit des transitions énergétiques discrètes ;
        \item il est caractéristique d'une espèce chimique ;
        \item il n'apparaît jamais en émission.
    \end{enumerate}
\end{enumerate}

\section{Vrai/faux exigeants}

\exostar{23}{4}
Justifier soigneusement.
\begin{enumerate}[label=\arabic*.,leftmargin=1cm]
    \item Quand une onde passe d'un milieu à un autre, sa période change.
    \item Un objet noir n'a pas de couleur.
    \item Une image virtuelle ne peut jamais être vue par un œil.
    \item Une lumière blanche est nécessairement monochromatique.
    \item Une radiation de plus courte longueur d'onde correspond à une plus grande fréquence.
    \item Si deux points sont séparés de $2\lambda$, ils vibrent en phase.
    \item La célérité d'une onde mécanique sur une corde dépend de la façon dont la source vibre.
    \item Deux photons de couleurs différentes ont nécessairement des énergies différentes.
\end{enumerate}

\section{Corrigé détaillé}

\begin{methodebox}
Pour exploiter ce corrigé intelligemment :
\begin{itemize}[leftmargin=1.2em]
    \item essaie vraiment avant de lire ;
    \item compare ta rédaction à la logique de la solution ;
    \item repère si ton erreur vient d'une formule, d'une conversion, ou du sens physique.
\end{itemize}
\end{methodebox}

\subsection*{Corrigé des QCM de démarrage}
\begin{corrbox}
Exercice 1 :
\begin{center}
\begin{tabular}{>{\bfseries}c cccccccccc}
\toprule
Question & 1 & 2 & 3 & 4 & 5 & 6 & 7 & 8 & 9 & 10\\
\midrule
Réponse & b & c & b & c & b & c & c & b & c & a\\
\bottomrule
\end{tabular}
\end{center}
\textbf{Commentaires utiles :} la fréquence est imposée par la source ; la célérité dépend du milieu ; dans le vide, la lumière vérifie $c=\lambda\nu$ ; plus $\nu$ augmente, plus l'énergie du photon augmente.
\end{corrbox}

\subsection*{Corrigé des vrai/faux}
\begin{corrbox}
Exercice 2 :
\begin{enumerate}[label=\arabic*.,leftmargin=0.8cm]
    \item \faux. Une onde transporte de l'énergie mais pas de matière à grande échelle.
    \item \faux. Une onde mécanique a besoin d'un milieu matériel ; l'amplitude n'y change rien.
    \item \vrai. La célérité dépend des propriétés du milieu.
    \item \vrai. La fréquence est fixée par la source qui vibre.
    \item \faux. En changeant de milieu, la fréquence reste la même ; ce sont la célérité et la longueur d'onde qui changent.
    \item \vrai. $\lambda=c/f$ dépend de $f$ donc de la source, et de $c$ donc du milieu.
    \item \faux. Une séparation de $\lambda/2$ correspond à une opposition de phase.
    \item \vrai. Pour un objet réel placé au-delà du foyer d'une lentille convergente, l'image réelle est renversée.
    \item \faux. Sous lumière bleue, un objet rouge absorbe le bleu et apparaît sombre, voire noir idéalisé.
    \item \vrai. Comme $E=hc/\lambda$, l'énergie est inversement proportionnelle à $\lambda$.
    \item \vrai. Les raies traduisent des transitions entre niveaux d'énergie discrets.
    \item \faux. Les deux modèles sont complémentaires selon le phénomène étudié.
\end{enumerate}
\end{corrbox}

\subsection*{Applications directes}
\begin{corrbox}
Exercice 3 :
\[
\tau=\frac{AB}{c}=\frac{1{,}5}{5{,}0}=\SI{0,30}{s}.
\]
Le point $B$ reproduit donc la perturbation observée en $A$ avec un retard de \SI{0,30}{s}.
\end{corrbox}

\begin{corrbox}
Exercice 4 :
\[
 d_{\text{total}}=c\Delta t = 340\times 0{,}80 = \SI{272}{m}.
\]
Le son effectue un aller-retour, donc la distance à l'obstacle vaut
\[
 d=\frac{272}{2}=\SI{136}{m}.
\]
On divise par 2 car l'écho correspond à un trajet aller puis retour.
\end{corrbox}

\begin{corrbox}
Exercice 5 :
\[
T=\frac{1}{f}=\frac{1}{25}=\SI{0,040}{s},\qquad \lambda = cT = 3{,}0\times 0{,}040 = \SI{0,12}{m}.
\]
Deux points en phase sont séparés de $\lambda=\SI{12}{cm}$, ou d'un multiple de cette valeur. En opposition de phase : $\lambda/2=\SI{6}{cm}$.
\end{corrbox}

\begin{corrbox}
Exercice 6 :
4 périodes en \SI{20}{ms} donnent
\[
T=\frac{20\ \si{ms}}{4}=\SI{5,0}{ms}=5{,}0\times10^{-3}\ \si{s}.
\]
Donc
\[
f=\frac{1}{T}=\frac{1}{5{,}0\times10^{-3}}=\SI{200}{Hz}.
\]
Si la période double, la fréquence est divisée par deux : le son devient plus grave.
\end{corrbox}

\begin{corrbox}
Exercice 7 : la distance entre deux crêtes successives est la longueur d'onde :
\[
\lambda = \SI{12}{cm}=\SI{0,12}{m}.
\]
Donc
\[
c=\frac{\lambda}{T}=\frac{0{,}12}{0{,}40}=\SI{0,30}{m.s^{-1}}.
\]
Si la fréquence double dans le même milieu, $c$ reste constant donc $\lambda = c/f$ est divisée par deux.
\end{corrbox}

\subsection*{Optique et couleurs}
\begin{corrbox}
Exercice 8 : avec $\overline{OA}=-\SI{30}{cm}$ et $f'=+\SI{10}{cm}$,
\[
\frac{1}{\overline{OA'}}-\frac{1}{\overline{OA}}=\frac{1}{f'}
\quad\Longrightarrow\quad
\frac{1}{\overline{OA'}}+\frac{1}{30}=\frac{1}{10}.
\]
Ainsi
\[
\frac{1}{\overline{OA'}}=\frac{1}{10}-\frac{1}{30}=\frac{2}{30}=\frac{1}{15}
\]
donc $\overline{OA'}=\SI{15}{cm}$.
Le grandissement vaut
\[
\gamma=\frac{\overline{OA'}}{\overline{OA}}=\frac{15}{-30}=-0{,}50.
\]
Donc
\[
A'B' = \gamma AB = -0{,}50\times \SI{2,0}{cm} = -\SI{1,0}{cm}.
\]
L'image mesure \SI{1,0}{cm}, elle est réelle et renversée.
\end{corrbox}

\begin{corrbox}
Exercice 9 :
\begin{enumerate}[label=\alph*)]
    \item Une image sur un écran est réelle, car les rayons lumineux y convergent effectivement.
    \item Si l'objet est entre le foyer et la lentille, l'image est virtuelle : on ne peut pas la projeter sur un écran.
    \item On peut déplacer la lentille par rapport au capteur/écran, ou modifier la focale.
    \item La mise au point permet de faire coïncider le plan image avec le capteur ou la rétine.
\end{enumerate}
\end{corrbox}

\begin{corrbox}
Exercice 10 :
\begin{enumerate}[label=\alph*)]
    \item Un objet cyan diffuse surtout le vert et le bleu, et absorbe le rouge.
    \item Sous lumière rouge, un objet jaune apparaît rouge car il diffuse le rouge.
    \item Sous lumière bleue, l'objet jaune absorbe le bleu et apparaît noir idéalisé.
    \item La synthèse additive concerne l'addition de lumières ; la synthèse soustractive concerne des pigments ou filtres qui retirent certaines composantes.
\end{enumerate}
\end{corrbox}

\subsection*{Photons et spectres}
\begin{corrbox}
Exercice 11 :
\[
E=\frac{hc}{\lambda}=\frac{6{,}63\times10^{-34}\times 3{,}00\times10^8}{650\times10^{-9}}
\approx 3{,}06\times10^{-19}\ \si{J}.
\]
Dans le visible, une longueur d'onde rouge comme \SI{650}{nm} correspond à un photon moins énergétique qu'un photon bleu ou violet.
\end{corrbox}

\begin{corrbox}
Exercice 12 :
\[
\nu=\frac{\Delta E}{h}=\frac{3{,}0\times10^{-19}}{6{,}63\times10^{-34}}\approx 4{,}52\times10^{14}\ \si{Hz}.
\]
Puis
\[
\lambda = \frac{c}{\nu}\approx \frac{3{,}00\times10^8}{4{,}52\times10^{14}}\approx 6{,}64\times10^{-7}\ \si{m}=\SI{664}{nm}.
\]
La radiation est dans le visible rouge. On obtient des raies car seules certaines transitions d'énergie sont possibles.
\end{corrbox}

\subsection*{Problèmes progressifs}
\begin{corrbox}
Exercice 13 : la lumière se propage à environ $3{,}0\times10^8\ \si{m.s^{-1}}$, donc son temps de parcours sur quelques kilomètres est négligeable à l'échelle de quelques secondes.
\[
d=340\times 5{,}4=\SI{1836}{m}\approx \SI{1,84}{km}.
\]
Le second observateur, \SI{680}{m} plus près de l'orage, entendra le tonnerre plus tôt de
\[
\Delta t = \frac{680}{340}=\SI{2,0}{s}.
\]
\end{corrbox}

\begin{corrbox}
Exercice 14 : 10 périodes en \SI{4,0}{s} donnent
\[
T=\frac{4{,}0}{10}=\SI{0,40}{s},\qquad f=\frac{1}{T}=\SI{2,5}{Hz}.
\]
Trois crêtes consécutives couvrent deux longueurs d'onde, donc
\[
2\lambda = \SI{45}{cm}\Rightarrow \lambda=\SI{22,5}{cm}=\SI{0,225}{m}.
\]
Alors
\[
c=\lambda f = 0{,}225\times 2{,}5 = \SI{0,5625}{m.s^{-1}}.
\]
Le retard à \SI{90}{cm} vaut
\[
\tau = \frac{0{,}90}{0{,}5625} \approx \SI{1,6}{s},
\]
donc le point commence à osciller après \SI{0,50}{s}.
\end{corrbox}

\begin{corrbox}
Exercice 15 : on a $f'=\SI{50}{mm}$ et $OA'=\SI{60}{mm}$.
\[
\frac{1}{OA'}-\frac{1}{OA}=\frac{1}{f'}
\Rightarrow \frac{1}{60}-\frac{1}{OA}=\frac{1}{50}.
\]
Donc
\[
-\frac{1}{OA}=\frac{1}{50}-\frac{1}{60}=\frac{1}{300}
\Rightarrow OA=-\SI{300}{mm}=-\SI{30}{cm}.
\]
Le grandissement vaut
\[
\gamma=\frac{60}{-300}=-0{,}20.
\]
La taille de l'image est donc
\[
A'B'=-0{,}20\times \SI{12}{cm}=-\SI{2,4}{cm}.
\]
L'image est renversée. Si la distance objet change, la position du plan image change aussi : il faut donc refaire la mise au point.
\end{corrbox}

\begin{corrbox}
Exercice 16 :
\begin{center}
\begin{tabular}{lcccc}
\toprule
Objet / éclairage & Blanc & Rouge & Vert & Bleu\\
\midrule
Objet rouge & rouge & rouge & noir & noir\\
Objet cyan & cyan & noir & vert & bleu\\
Objet blanc & blanc & rouge & vert & bleu\\
\bottomrule
\end{tabular}
\end{center}
Un objet blanc diffuse toutes les composantes reçues ; un objet rouge diffuse surtout le rouge ; un objet cyan diffuse le vert et le bleu. Oui, un objet cyan peut apparaître noir sous lumière rouge, car il absorbe cette composante.
\end{corrbox}

\begin{corrbox}
Exercice 17 : plus $\lambda$ est petite, plus l'énergie est grande. Donc : \SI{410}{nm} est la plus énergétique, puis \SI{434}{nm}, puis \SI{486}{nm}. La raie de \SI{820}{nm} serait dans l'infrarouge proche, hors du visible. La transition associée à \SI{410}{nm} est celle de plus grande variation d'énergie.
\end{corrbox}

\subsection*{Problèmes plus exigeants}
\begin{corrbox}
Exercice 18 :
\[
\Delta t = \SI{4,0}{ms}=4{,}0\times10^{-3}\ \si{s},\qquad \Delta d = c\Delta t = 340\times 4{,}0\times10^{-3}=\SI{1,36}{m}.
\]
La source est donc plus proche de $M_1$ que de $M_2$ de \SI{1,36}{m}. Si elle est entre les deux microphones, en notant $x=M_1S$, on a
\[
(3{,}40-x)-x=1{,}36 \Rightarrow 3{,}40-2x=1{,}36 \Rightarrow x=\SI{1,02}{m}.
\]
Si elle est à gauche de $M_1$, les distances aux micros diffèrent toujours de \SI{3,40}{m}, ce qui ne convient pas : cette possibilité est exclue pour cette géométrie linéaire. En dimension 2 ou 3, un seul retard ne suffit généralement pas ; il faut plusieurs capteurs.
\end{corrbox}

\begin{corrbox}
Exercice 19 : on cherche $f'$ avec $OA=-\SI{2,0}{m}=-\SI{2000}{mm}$ et $OA'=\SI{52}{mm}$.
\[
\frac{1}{f'}=\frac{1}{52}-\frac{1}{2000}\approx 0{,}01873
\Rightarrow f'\approx \SI{53,4}{mm}.
\]
Une focale plus grande sans déplacer le capteur modifie la position de l'image : elle n'est plus nette. Une focale plus grande augmente en général le grandissement, ce qui correspond à l'idée d'un zoom optique réel : on change la formation géométrique de l'image, pas seulement son affichage numérique.
\end{corrbox}

\begin{corrbox}
Exercice 20 : avec $OA=-\SI{1,20}{m}=-\SI{1200}{mm}$ et $f'=\SI{100}{mm}$,
\[
\frac{1}{OA'}-\frac{1}{-1200}=\frac{1}{100}
\Rightarrow \frac{1}{OA'}=\frac{1}{100}-\frac{1}{1200}=\frac{11}{1200}.
\]
Donc
\[
OA'\approx \SI{109}{mm}.
\]
Le grandissement vaut
\[
\gamma=\frac{109}{-1200}\approx -0{,}091.
\]
Pour un objet de \SI{3,0}{mm}, l'image mesure environ
\[
|A'B'|\approx 0{,}091\times 3{,}0 \approx \SI{0,27}{mm}.
\]
L'énergie du photon de \SI{486}{nm} vaut
\[
E=\frac{hc}{\lambda}\approx \frac{6{,}63\times10^{-34}\times3{,}00\times10^8}{486\times10^{-9}}\approx 4{,}09\times10^{-19}\ \si{J}.
\]
Une raie spectrale identifie une espèce car les niveaux d'énergie sont spécifiques. On combine ici la formation d'image par lentille et la nature quantifiée des échanges lumière-matière. Erreurs fréquentes : oublier de convertir les unités ; confondre taille algébrique et taille géométrique positive.
\end{corrbox}

\begin{corrbox}
Exercice 21 : connaître seulement $T$ ne permet ni de déduire $\lambda$ ni de déduire $c$. Il faut au moins une mesure spatiale (par exemple $\lambda$) ou une mesure de retard/célérité. Deux protocoles possibles :
\begin{itemize}[leftmargin=1.2em]
    \item mesurer la distance entre deux crêtes successives pour obtenir $\lambda$, puis calculer $c=\lambda/T$ ;
    \item mesurer le temps mis par une perturbation pour parcourir une distance connue, puis calculer $c=d/\Delta t$.
\end{itemize}
Cet exercice teste la compréhension des dépendances entre grandeurs : aucune formule ne peut créer une information absente.
\end{corrbox}

\subsection*{Corrigé des QCM d'approfondissement et vrai/faux exigeants}
\begin{corrbox}
Exercice 22 :
\begin{center}
\begin{tabular}{>{\bfseries}c l}
\toprule
Question & Réponse(s) correcte(s)\\
\midrule
1 & b\\
2 & b\\
3 & a et c\\
4 & b et c\\
\bottomrule
\end{tabular}
\end{center}
\end{corrbox}

\begin{corrbox}
Exercice 23 :
\begin{enumerate}[label=\arabic*.,leftmargin=0.8cm]
    \item \faux. La fréquence et donc la période sont imposées par la source ; elles restent inchangées lors d'un changement de milieu.
    \item \faux. Un objet noir absorbe l'essentiel de la lumière reçue ; ce n'est pas une absence de nature physique mais un comportement optique simplifié.
    \item \faux. Une image virtuelle peut être vue à travers une lentille, même si elle n'est pas projetable sur un écran.
    \item \faux. La lumière blanche est polychromatique.
    \item \vrai. Comme $c=\lambda\nu$, à célérité fixée dans le vide, plus $\lambda$ est petite, plus $\nu$ est grande.
    \item \vrai. Toute séparation égale à un multiple entier de $\lambda$ correspond à des points en phase.
    \item \faux. La célérité dépend du milieu ; la source fixe surtout la fréquence.
    \item \vrai. Des couleurs différentes correspondent à des fréquences différentes, donc à des énergies différentes.
\end{enumerate}
\end{corrbox}

\begin{ouverturebox}
Pour la Terminale, il sera très utile d'être parfaitement à l'aise avec :
\begin{itemize}[leftmargin=1.2em]
    \item la distinction entre \textbf{ce qui dépend de la source} et \textbf{ce qui dépend du milieu} ;
    \item la lecture qualitative et quantitative des \textbf{spectres} ;
    \item l'idée qu'un même phénomène lumineux peut demander plusieurs \textbf{modèles complémentaires}.
\end{itemize}
\end{ouverturebox}

\end{document}
